\section{Introduction}\label{sec:introduction}

A homogeneous ideal is constrained both by the equations it contains and by the dimensions of its graded pieces. The Eisenbud--Green--Harris conjecture predicts a particularly concrete form for this constraint: a regular sequence can be replaced by the corresponding pure powers without changing the Hilbert function of a containing ideal. We prove this assertion over the complex numbers in its Artinian formulation, for arbitrary degrees, and then deduce the full characteristic-zero statement for regular sequences of arbitrary positive length. The main construction provides more than the Hilbert-function comparison: it replaces the polynomial variables by commuting linear forms with division-algebra coefficients and retains an ordered basis of the entire polynomial algebra.

Let $S=\C[x_1,\ldots,x_n]$ with its standard grading. For a homogeneous ideal $I\subseteq S$, write
\[
 H_{S/I}(d)=\dim_\C(S/I)_d\qquad(d\ge0).
\]
A sequence $f_1,\ldots,f_n$ of positive-degree homogeneous elements is \emph{regular} if multiplication by $f_i$ is injective on $S/(f_1,\ldots,f_{i-1})$ for each $i$. Its quotient has Hilbert series
\[
 \Hilb_{S/(f_1,\ldots,f_n)}(t)
 =\prod_{i=1}^n(1+t+\cdots+t^{a_i-1}),\qquad a_i=\deg f_i,
\]
so every ideal containing it has an Artinian quotient.

Eisenbud, Green, and Harris arrived at the conjecture through extensions
of Castelnuovo theory. Section~4 of \cite{EGH93} formulates the quadratic
case and then proposes the arbitrary-degree pure-power comparison.
Their Cayley--Bacharach interpretation already appears in
\cite[Section~3]{EGH93} and is developed further in
\cite[Section~2.1]{EGH96}. The pure-power model turns the geometric
constraints into a problem about monomials with bounded exponents.
The construction below supplies that model for an arbitrary homogeneous
regular sequence.

We first state this construction. If $k/\C$ is a field extension and
$\Delta$ is a central division $k$-algebra, then
$\Delta[x_1,\ldots,x_n]$ denotes the polynomial algebra with central
variables $x_j$. The coefficients in $\Delta$ have degree zero. A
collection of its linear forms may commute even when their individual
coefficients do not.

\begin{theorem}[Commuting-form construction]\label{thm:construction}
Let $n\ge1$, let $2\le a_1\le\cdots\le a_n$, and let
$f_1,\ldots,f_n$ be a homogeneous regular sequence in
$\C[x_1,\ldots,x_n]$ with $\deg f_i=a_i$. There exist a finitely
generated field extension $k/\C$, a finite-dimensional central division
$k$-algebra $\Delta$, and a commutative graded $k$-subalgebra
\[
 k[x_1,\ldots,x_n]\subseteq A\subseteq\Delta[x_1,\ldots,x_n]
\]
generated by finitely many pairwise commuting linear forms, with
distinguished $t_1,\ldots,t_n\in A_1$, such that
\begin{equation}\label{eq:construction-triangular}
 t_i^{a_i}=c_i f_i+\sum_{l<i}B_{il}f_l+\sum_{d>i}C_{id}t_d
 \qquad(1\le i\le n).
\end{equation}
Here $c_i\in k^\times$, $B_{il}\in A_{a_i-a_l}$, and
$C_{id}\in A_{a_i-1}$. Moreover, $A$ is finite over $k[x]$, and
\[
 \{t_1^{\alpha_1}\cdots t_n^{\alpha_n}:
             \alpha\in\Z_{\ge0}^n\}
\]
is a basis of the full polynomial algebra $\Delta[x]$ as a left
$\Delta$-vector space.
\end{theorem}

The basis is not restricted to the complete-intersection quotient. In
particular, it describes every graded piece before any additional ideal
is imposed. This is the interface used for resolution transfer in the
companion paper \cite[Lemma~4.1]{BettiCompanion}.
For Hilbert functions, a least-exponent argument already gives the
following consequence.

\begin{corollary}[Artinian Eisenbud--Green--Harris]\label{thm:main}
Let $n\ge1$ and $2\le a_1\le\cdots\le a_n$. Suppose that a homogeneous ideal $I\subseteq\C[x_1,\ldots,x_n]$ contains a regular sequence $f_1,\ldots,f_n$ with $\deg f_i=a_i$. There is a monomial ideal $J\subseteq\C[x_1,\ldots,x_n]$ such that
\[
 x_i^{a_i}\in J\quad(1\le i\le n),
 \qquad
 H_{S/J}(d)=H_{S/I}(d)\quad\text{for every }d\ge0.
\]
\end{corollary}

Thus the Artinian Eisenbud--Green--Harris conjecture is resolved positively over $\C$. One ideal $J$ realizes the entire Hilbert function, and there is no restriction on the further generators of $I$.

To describe the lex-plus-powers formulation, order variables by
$x_1>\cdots>x_n$, and let $J$ be the monomial ideal in the corollary.
In degree $d$, among monomials outside
$P=(x_1^{a_1},\ldots,x_n^{a_n})$, take the
$q_d=\dim_\C J_d-\dim_\C P_d$ lexicographically largest ones. The
Clements--Lindstr\"om theorem says that these degreewise choices, together
with $P$, form an ideal whenever the Hilbert function comes from an ideal
containing $P$ \cite{CL69}. Thus Corollary~\ref{thm:main} also gives one
lex-plus-powers ideal with the entire prescribed Hilbert function. We use
the exact compression statement in
\cite[Theorem~2.1(i)]{CDS22}; it is not a premise of
Theorem~\ref{thm:construction}.

The Hilbert-function conclusion is not restricted to Artinian quotients
or to the coefficient field \(\C\).

\begin{corollary}[Eisenbud--Green--Harris in characteristic zero]
\label{cor:characteristic-zero}
Let \(F\) be a field of characteristic zero, let
\(S_F=F[x_1,\ldots,x_n]\) be standard graded, and let
\[
 1\le c\le n,\qquad 2\le a_1\le\cdots\le a_c.
\]
If a homogeneous ideal \(I\subseteq S_F\) contains a homogeneous regular
sequence \(f_1,\ldots,f_c\) with \(\deg f_i=a_i\), there is a monomial ideal
\(J\subseteq S_F\) containing
\(P_F=(x_1^{a_1},\ldots,x_c^{a_c})\) such that
\[
 \dim_F(S_F/J)_d=\dim_F(S_F/I)_d\qquad(d\ge0).
\]
Moreover, \(J\) may be chosen lex-plus-powers: in each degree its monomials
outside \(P_F\) are the lexicographically greatest
\(\dim_F I_d-\dim_F(P_F)_d\) monomials outside \(P_F\), for the order
\(x_1>\cdots>x_n\).
\end{corollary}

The proof in Section~\ref{subsec:generality} combines the Artinian result
with a reduction to full regular sequences and descent through a finitely
generated coefficient field. The division-algebra construction of
Theorem~\ref{thm:construction} itself concerns full regular sequences over
\(\C\). Section~\ref{sec:consequences} also gives local-cohomology
inequalities for the resulting lex-plus-powers ideal and a
scheme-theoretic quadratic Cayley--Bacharach bound.

\subsection{Earlier results on EGH}

Results depending only on the degrees of the regular sequence have
covered important ranges. Caviglia and Maclagan proved EGH when
$a_i>\sum_{j<i}(a_j-1)$ for $i\ge3$ \cite{CM08}; Caviglia and De Stefani
extended this to the non-strict inequalities
\cite[Theorem~A]{CDeS20}. For quadratic sequences, Chen proved the case
of at most four variables \cite[Theorem~2.2]{Chen12}. Caviglia,
De Stefani, and Sbarra give an independent proof for five quadrics,
alongside a systematic account of the conjecture
\cite[Theorem~4.14]{CDS22}.

Other results use additional structure of the forms or the containing
ideal. Abedelfatah proves EGH when all forms of the regular sequence
split into linear factors \cite[Corollary~4.3]{Abedelfatah15}.
Cooper treats subsets of finite reduced complete intersections of
types $(a,b)$, $(2,a,b)$, and $(3,a,a)$ \cite{Cooper12}.
Chong uses liaison to treat sequentially bounded licci ideals with a
minimal first link and a minimally contained regular sequence; his
results include height-three Gorenstein ideals under the corresponding
minimal-containment condition
\cite[Theorem~3 and Corollary~11]{Chong16}.

Recent work addresses both small-degree cases and asymptotic ranges.
Abedelfatah proves EGH for quadratic almost complete intersections in
six variables \cite[Theorem~4.3]{Abedelfatah26}. Kuzmanovski obtains
Hilbert-function comparisons through prescribed degree ranges when the
initial degree is sufficiently large
\cite[Theorems~1.4--1.5]{Kuzmanovski26}. The construction here applies to
arbitrary regular sequences of the stated degrees and gives one ideal
with the same Hilbert function in every degree.

\subsection{Proof strategy and its methods}

The triangular identities in Theorem~\ref{thm:construction} make the
\(t_i\) a regular system of parameters on the finite module
\(\Delta[x]\) over \(A\). Since they have degree one, a Hilbert-series
count gives the full ordered basis. Least exponents of elements of an
extended ideal then form one upward-closed set, defining the ordinary
monomial ideal. Section~\ref{sec:reduction} establishes this algebraic
reduction before the construction begins.

We construct the rows from \(i=n\) down to \(i=1\). At a stage with
\(b=a_i\), the aim is to express a suitable right-hand side of the new
row as a sum \(\ell_1^b+\cdots+\ell_M^b\) of powers of existing forms,
then combine these powers two at a time until the sum becomes \(t_i^b\).
Two difficulties govern the construction: the binary addition matrices
must have coefficients in a division algebra containing the old
coefficients, and the earlier power identities must remain available
when the next degree is smaller.

To retain those identities, we use a finite-dimensional space \(U\)
of formal linear forms containing the variables \(x_j\) and the forms
\(t_d\) already constructed, together with an ideal
\(Q\subseteq\Sym U\) generated in degrees one and two. Evaluation
recovers the actual commuting forms, and every earlier identity holds
modulo \(Q\). Let \(\mathfrak j_b\) be the degree-\(b\) part of
\[
 (Q,f_1,\ldots,f_i,t_{i+1},\ldots,t_n)\subseteq\Sym U.
\]
A linear functional \(\lambda\) on \(\Sym^b U\) that annihilates
\(\mathfrak j_b\) defines the
degree-\(b\) polynomial \(w_\lambda(\ell)=\lambda(\ell^b)\) on \(U\).
If \(W\) is their space, then
\[
 \sum_{m=1}^M\ell_m^b\in\mathfrak j_b
 \quad\Longleftrightarrow\quad
 \sum_{m=1}^M w(\ell_m)=0\qquad(w\in W).
\]
Let \(X\subseteq\PP(U^{\oplus M})\) be the projectivization of the
common zero set of these equations.

The geometric safeguard is that the projective zero set of \(Q\)
contains no curve of bounded small degree. Section~\ref{sec:low-relations}
constructs such data from independent quadratic equations, proves that
finite homogeneous enlargement preserves the property, and uses it to
show that every nonzero \(w\in W\) has an irreducible factor of
multiplicity one. For sufficiently many blocks, Section~\ref{sec:parameters}
then proves that \(X\) is an integral
complete intersection and that the nonzero cone over its smooth locus
has arbitrarily high connectivity. The argument allows nonisolated
singularities and requires only one simple factor in each member of
\(W\).

The binary addition matrices come from a curve with three sections
satisfying \(z^b=u^b+v^b\). A generic degree-zero line bundle makes
multiplication by these sections linear in \(u,v\). This use of matrix
pencils belongs to the generalized Clifford-algebra and vector-bundle
correspondence developed by van den Bergh \cite{VanDenBergh87} and
Kulkarni \cite{Kulkarni03}; the trivial finite-pushforward formulation
appears in \cite[Proposition~2.5]{CKM12}.
Section~\ref{sec:curves} constructs the required curve and pencils
explicitly, together with a finite group action. A matrix representation
alone does not make the coefficient algebra division.

The two Picard varieties of the curve serve different purposes:
degree-zero line bundles parametrize the matrices, while the
degree-one Picard variety supplies a free group action and maps carrying
prescribed characters. Equivariant descent first gives a central simple
algebra. To prove that it is division, we show that the dimension of
every finite right module is divisible by the dimension of the algebra,
both measured over its center. A nonzero proper right ideal would
contradict this divisibility.

Section~\ref{sec:descent} uses a determinant-of-cohomology
correction to relate module dimensions to Euler characteristics on a
product of Picard varieties. The free action also gives divisibility
of Euler characteristics for the corresponding family over \(X\),
after twisting by \(\cO_X(l)\). The correction uses the formalism of
Knudsen and Mumford
\cite{KnudsenMumford76}, and its numerical effect follows from
Grothendieck's Riemann--Roch theorem \cite{BorelSerre58}.
Section~\ref{sec:index} uses the character maps, connectivity of the
smooth cone, and integral topological \(K\)-theory to compare these
Euler characteristics and prove that the algebra is division while
retaining the old coefficient algebra.

The use of highly connected algebraic parameter spaces and finite
quotients is related to the Godeaux--Serre approximation method
\cite[Section~5]{Totaro99}. Topological restrictions on an algebraic
index have a close precedent in
Antieau and Williams \cite[Theorem~3.4]{AntieauWilliams13}.
Here the determinant correction and integral \(K\)-theory pushforward
separate divisibility by the new group order from divisibility by the
old algebra's degree. Consequently the required connectivity can be
fixed before enlarging the number of summands. Finally,
Section~\ref{sec:assembly} verifies the characters that make the new
output invariant and retains the new row through a finite extension
defined by linear and quadratic relations. This closes the induction.

When every $a_i=2$, Corollary~\ref{thm:main} gives the quadratic assertion.
The complete-intersection Hilbert series is then $(1+t)^n$, and the
standard monomials of a monomial ideal containing $x_1^2,\ldots,x_n^2$
correspond to a family of subsets of $\{1,\ldots,n\}$ closed under
taking subsets. Here this
special case follows from the same ordered-basis construction as all
other degree sequences; no separate quadratic theorem is required.

Figure~\ref{fig:construction} records the dependencies within one induction
step and the return to its hypotheses before the next row.

\begin{figure}[!htbp]
\centering
\begin{tikzpicture}[
 box/.style={draw=linkblue!70,rounded corners=2pt,fill=linkblue!3,align=center,text width=38mm,minimum height=12mm,font=\small,inner sep=4pt},
 arrow/.style={-{Stealth[length=2mm]},semithick,draw=linkblue!80},
 every node/.style={font=\small}]
\node[box] (low) at (0,2.1) {Low-degree relations\\and no small curves};
\node[box] (param) at (5,2.1) {Constraint space $X$\\with highly connected\\punctured smooth cone};
\node[box] (curves) at (0,0) {Binary addition\\matrices on curves};
\node[box] (division) at (5,0) {Division algebra\\extension};
\node[box] (row) at (10,0) {New triangular\\power identity};
\node[box] (ideal) at (10,-2.1) {Ordered basis and\\one monomial ideal};
\draw[arrow] (low)--(param);
\draw[arrow] (param)--(division);
\draw[arrow] (curves)--(division);
\draw[arrow] (division)--(row);
\draw[arrow] (row)-- node[right,align=left]{after all\\rows} (ideal);
\draw[arrow] (row.north) -- (10,3.2) -- node[above]{finite extension; retain low-degree relations} (0,3.2) -- (low.north);
\end{tikzpicture}
\caption{The induction separates the geometry of the constraints from the binary addition matrices. The return arrow preserves the hypotheses for the next row; monomial extraction is performed only after the construction is complete.}
\label{fig:construction}
\end{figure}

\FloatBarrier
\subsection{Conventions and standard tools}\label{subsec:conventions}

Every center $k$ used in the construction is a finitely generated field extension of the original $\C$. The degree of a central simple algebra is the square root of its dimension over its center. Linear forms with division-algebra coefficients always commute as whole elements of the polynomial algebra; their individual coefficients need not commute. When new tensor factors are used, we will specify the stronger coefficientwise commutativity available between distinct factors.

Each such $k$ admits an abstract embedding into $\C$: its algebraic closure has the same transcendence degree over $\Q$ as $\C$. These embeddings need not fix the original constants. They are used only for complex topology and numerical base-change comparisons, after the relevant geometric hypotheses have been verified. All actual algebra extensions continue to contain the original $\C$ centrally. In particular, division is proved over the unembedded field by rank divisibility, never inferred from a complex specialization. Embeddings chosen at successive stages need not extend one another.

Projective spaces parameterize lines. For a vector space $U$ of formal forms, $\Sym U$ gives those forms degree one, and its geometric projective zero sets lie in $\PP(U^\vee)$. Polynomial functions $w$ on $U$, in contrast, are evaluated on the forms $\ell\in U$. We distinguish these two roles when using differentiation.

We use standard coherent-sheaf and topological $K$-theory with their usual complex orientations. Pullbacks of coherent classes are derived whenever ordinary pullback is not exact. On smooth quasi-projective varieties, finite locally free resolutions give the corresponding topological $K$-classes. For smooth projective varieties, the Riemann--Roch formulas are
\[
 \chi(V,\alpha)=\int_V\ch(\alpha)\td(V),
 \qquad
 \ch(Rp_*\alpha)=p_*\bigl(\ch(\alpha)\td(T_p)\bigr)
\]
for projection $p$ off a fixed smooth projective factor. These formulas apply to coherent classes, without assuming that their individual higher direct images are locally free; see \cite[Sections~4, 6, and~7]{BorelSerre58}. The singular-space Euler-characteristic argument needed below is proved separately. The integral projective-space basis and complex Thom class in topological $K$-theory are recalled at their point of use from \cite{Hatcher}.
